\documentclass[11pt]{article}
\usepackage[margin=1in]{geometry}
\usepackage[T1]{fontenc}
\usepackage{graphicx}
\usepackage{hyperref}
\usepackage{enumitem}
\setlength{\parskip}{0.5em}
\setlength{\parindent}{0pt}

\title{diffusion models for image generation}
\author{arXiv GraphRAG --- phase-6 deep research}
\date{\today}

\begin{document}
\maketitle

\begin{abstract}
The review is excellent in its historical lineage but lacks coverage of recent preference alignment (DPO) and the specific shift toward one-step consistency models.
\end{abstract}

\section{From Noise to Nuance: A Comprehensive Review of Diffusion Models for Image Generation}

The ability to synthesize high-fidelity, diverse, and semantically coherent imagery has long been a cornerstone of generative modeling research. While previous paradigms struggled to capture the intricate textures and global structures of the natural world, the emergence of diffusion-based frameworks has redefined the boundaries of machine-driven creativity. This review begins by contextualizing these advancements within the broader foundations of generative modeling and early diffusion, tracing the mathematical evolution from simple stochastic processes to the sophisticated mechanisms that allow for the reconstruction of structured data from pure noise.

The trajectory of the field is marked by a significant pivot following the rise of denoising diffusion probabilistic models, which established a superior framework for sample quality and stability. This technological leap paved the way for the integration of semantic guidance, specifically through latent space and text-conditional diffusion, which enabled the alignment of visual synthesis with complex linguistic prompts. As these models grew in complexity, the research community shifted its focus toward practical deployment, driving innovations in accelerating sampling and efficient diffusion, as well as specialized techniques for image restoration and conditional control to enhance user agency and output precision.

As the capabilities of these models mature, the scope of application is rapidly moving beyond static two-dimensional pixels. The current landscape is characterized by expanding dimensions, specifically through the development of video and 3D generation techniques that attempt to maintain temporal and spatial consistency. Finally, this review explores the theoretical frontier, where the focus shifts from empirical scaling to the fundamental study of manifold geometry and distributional modeling, seeking to provide a more rigorous mathematical understanding of how diffusion processes navigate the high-dimensional spaces of visual data.

\subsection{Foundations of Generative Modeling and Early Diffusion}

The evolution of generative modeling has transitioned from establishing fundamental probabilistic frameworks to the development of highly efficient diffusion-based architectures. Early advancements focused on optimizing latent variable models and adversarial training. Variational Autoencoders (VAEs) introduced the reparameterization trick to enable stochastic gradient optimization of the variational lower bound, facilitating efficient learning in directed probabilistic models with continuous latent variables [\texttt{arXiv:1312.6114}]. In contrast, Generative Adversarial Networks (GANs) bypassed Markov chains and unrolled approximate inference by framing generation as a minimax two-player game between a generator and a discriminator [\texttt{arXiv:1406.2661}]. While GANs provided a path to capture data distributions without complex inference, they faced stability challenges. This was partially addressed by the Two Time-Scale Update Rule (TTUR), which utilized individual learning rates for the generator and discriminator to ensure convergence to a local Nash equilibrium [\texttt{arXiv:1706.08500}]. Subsequent architectural innovations, such as style-based generators, improved GANs by enabling the unsupervised separation of high-level attributes from stochastic variations through scale-specific control [\texttt{arXiv:1812.04948}].

A parallel trajectory emerged from non-equilibrium statistical physics, forming the bedrock of modern diffusion models. Early work established the concept of systematically destroying data structure through an iterative forward diffusion process to learn a reverse generative process [\texttt{arXiv:1503.03585}]. This concept was further refined through score-based modeling, which utilized Langevin dynamics to estimate the gradients of the data distribution across multiple levels of Gaussian noise perturbation [\texttt{arXiv:1907.05600}].

The field reached a critical juncture with the unification of these approaches. Denoising Diffusion Probabilistic Models (DDPMs) established a novel connection between diffusion processes and denoising score matching, utilizing a weighted variational bound to achieve high-quality image synthesis [\texttt{arXiv:2006.11239}]. While DDPMs achieved state-of-the-art performance, such as a FID of 3.17 on CIFAR10, they suffered from uncompetitive log-likelihoods [\texttt{arXiv:2006.11239}]. This was addressed by framing generative modeling through Stochastic Differential Equations (SDEs), which unified score-based and diffusion models into a single framework, enabling exact likelihood computation and new sampling procedures [\texttt{arXiv:2011.13456}].

Despite these mathematical advancements, the slow sampling speed of diffusion models remained a significant bottleneck. Denoising Diffusion Implicit Models (DDIMs) addressed this by constructing non-Markovian diffusion processes that share the same training objective as DDPMs but allow for much faster, deterministic sampling [\texttt{arXiv:2010.02502}]. DDIMs achieved a 10× to 50× speedup in wall-clock time while maintaining high sample quality (e.g., a CIFAR10 FID of 4.16 at 100 steps), though they inherently trade off sample quality for computational speed [\texttt{arXiv:2010.02502}]. Together, these developments transitioned generative modeling from the adversarial competition of GANs toward the structured, gradient-based refinement characteristic of modern diffusion frameworks.

\subsection{The Rise of Denoising Diffusion Probabilistic Models}

The evolution of diffusion models has been characterized by a transition from foundational score-based modeling to highly efficient, high-resolution generative frameworks. Early advancements focused on the stability and scalability of score-based generative models. For instance, \texttt{arXiv:2006.09011} introduced theoretical analyses for learning in high-dimensional spaces and utilized an exponential moving average (EMA) of model weights to enhance training stability, successfully scaling models from 64x64 to 256x256 resolutions. This emphasis on stability and scaling was further refined in the context of likelihood optimization; \texttt{arXiv:2107.00630} proposed Variational Diffusion Models (VDM), which optimized noise schedules jointly with the model using a simplified Variational Lower Bound (VLB) based on signal-to-noise ratio (SNR) to achieve state-of-the-art density estimation and compression.

A significant bottleneck in early diffusion frameworks was the slow inference speed required by Markovian sampling processes. \texttt{arXiv:2010.00713} addressed this by introducing Denoising Diffusion Implicit Models (DDIMs), which constructed non-Markovian diffusion processes that share the same training objective as standard DDPMs. This innovation allowed for deterministic, implicit modeling and provided a substantial 10× to 50× speedup in wall-clock time while maintaining high sample quality, such as a CIFAR10 FID of 4.16 at 100 steps. However, this acceleration involves an inherent trade-off, as sample quality may decrease alongside increased computational speed. Complementing these sampling improvements, \texttt{arXiv:2102.09672} improved the underlying probabilistic models by learning the variances of the reverse diffusion process, thereby achieving more competitive log-likelihoods and reducing the required number of forward passes.

As the field moved toward complex data representations, researchers explored specialized architectures and latent spaces. \texttt{arXiv:2106.15282} utilized a cascaded approach, employing a pipeline of multiple diffusion models and conditioning augmentation to achieve high-fidelity class-conditional images, reaching an FID of 4.88 at 256x256. In the realm of text-to-image synthesis, \texttt{arXiv:2111.14822} moved the diffusion process into a vector quantized (VQ-VAE) latent space. By implementing a mask-and-replace strategy, this model avoided error accumulation and achieved generation speeds 15x faster than traditional autoregressive methods. Beyond pure image synthesis, the principles of gradient estimation in diffusion were extended to other modalities, such as WaveGrad (\texttt{arXiv:2009.00713}), which applied score-based modeling to high-fidelity waveform generation conditioned on mel-spectrograms. Finally, the integration of natural language supervision, exemplified by the contrastive pre-training in \texttt{arXiv:2103.00020}, provided a foundation for the zero-shot transfer capabilities that increasingly inform modern generative pipelines.

\subsection{Latent Space and Text-Conditional Diffusion}

The evolution of diffusion models has been defined by a shift from pixel-space operations to efficient latent-space manipulations and increasingly sophisticated text-conditioning mechanisms. A foundational advancement in this trajectory is the Latent Diffusion Model (LDM), which applies denoising processes within the latent space of pretrained autoencoders rather than the high-dimensional pixel space \texttt{arXiv:2112.10752}. This approach significantly reduces the computational requirements for both training and inference while enabling high-resolution synthesis \texttt{arXiv:2112.10752}. The scalability of this paradigm is further demonstrated by SDXL, which extends the LDM framework by utilizing a UNet backbone three times larger than its predecessor, incorporating a second text encoder, and employing multiple conditioning schemes to enhance visual fidelity \texttt{arXiv:2307.01952}.

Text-conditioning has transitioned from simple guidance to deep semantic integration. Early efforts like GLIDE explored the efficacy of CLIP guidance versus classifier-free guidance for photorealistic generation and editing \texttt{arXiv:2112.10741}. The importance of guidance mechanisms was formalized by the introduction of classifier-free guidance, which jointly trains conditional and unconditional models to allow a tunable trade-off between sample fidelity and mode coverage without the need for an external classifier \texttt{arXiv:2207.12598}.

A critical tension exists between the complexity of the diffusion backbone and the complexity of the text encoder. While architectures like Diffusion Transformers (DiT) emphasize scalability by replacing the traditional U-Net with a transformer architecture operating on latent patches \texttt{arXiv:2212.09748}, other research suggests that the bottleneck for alignment lies in language understanding. Specifically, it has been observed that increasing the size of large, generic, text-only pretrained language models (such as T5) boosts image-text alignment and fidelity more effectively than simply increasing the size of the diffusion model itself \texttt{arXiv:2205.11487}. This is complemented by hierarchical approaches, such as those using CLIP latents, which introduce an intermediate step by generating CLIP image embeddings as a prior to improve diversity and facilitate language-guided manipulation \texttt{arXiv:2204.06125}.

Finally, the ability to impose precise spatial constraints on these text-conditional models has been addressed through architectural modularity. ControlNet provides a method for adding spatial conditioning (e.g., depth or human pose) to locked, pretrained text-to-image models using zero-initialized convolution layers \texttt{arXiv:2302.05543}. This allows for the injection of structural control without disrupting the learned semantic knowledge of the original model, effectively bridging the gap between high-level text descriptions and low-level spatial precision.

\subsection{Accelerating Sampling and Efficient Diffusion}

The high computational cost of diffusion models stems from the iterative nature of the sampling process, prompting research into both mathematical solver optimization and structural modifications to the diffusion trajectory. A primary approach involves treating the sampling process as an Ordinary Differential Equation (ODE) to be solved more efficiently. DPM-Solver (\texttt{arXiv:2206.00927}) provides an exact formulation for diffusion ODEs by analytically computing the linear part of the solution via an exponentially weighted integral, achieving a 4-16x speedup and high-quality generation in as few as 10-20 steps (e.g., an FID of 2.87 in 20 steps on CIFAR10) (\texttt{arXiv:2206.00927}). This is extended by DPM-Solver++ (\texttt{arXiv:2211.01095}), which introduces a high-order solver and a multistep variant specifically designed to address instability issues during guided sampling in both pixel and latent spaces (\texttt{arXiv:2211.01095}).

While these solvers optimize the existing sampling trajectory, other methods attempt to reshape the trajectory itself to be more efficient. Rectified Flow (\texttt{arXiv:2209.03003}) utilizes nonlinear least squares optimization to learn straight paths between distributions, providing a unified framework for generative modeling and domain transfer (\texttt{arXiv:2209.03003}). Similarly, Hierarchical Koopman Diffusion (\texttt{arXiv:2510.12220}) seeks to achieve one-step sampling by lifting nonlinear dynamics into a latent space governed by globally linear operators through Koopman operator theory (\texttt{arXiv:2510.12220}).

A second research direction focuses on reducing the number of required evaluations through model distillation or structural decomposition. Progressive Distillation (\texttt{arXiv:2202.00512}) employs an iterative procedure to halve the required sampling steps, utilizing new parameterizations to maintain stability during few-step sampling (e.g., achieving an FID of 3.0 in 4 steps on CIFAR-10) (\texttt{arXiv:2202.00512}). In contrast to this iterative reduction, MSF (\texttt{arXiv:2501.13349}) achieves a 4x speedup by decomposing the denoising target into a low-frequency base signal and a high-frequency residual signal using multi-scale latent factorization (\texttt{arXiv:2501.13349}). This structural approach to efficiency is also seen in Accelerated Image-Aware Generative Diffusion (\texttt{arXiv:2408.08306}), which reduces the sampling requirement from \textasciitilde{}1000 steps to 200-500 by implementing an exponentially time-decaying Signal to Noise Ratio (SNR) (\texttt{arXiv:2408.08306}).

Finally, while much of the literature prioritizes speed, some work emphasizes the quality of the denoising process. For instance, redefining the diffusion process to match specific noise models—such as Gamma or Poisson—can prevent the over-smoothing of results and better preserve image details compared to standard Gaussian-based approaches (\texttt{arXiv:2302.02398}).

\subsection{Image Restoration and Conditional Control}

The evolution of diffusion models has transitioned from general generative priors to specialized frameworks capable of precise image-to-image translation and high-fidelity restoration. Early advancements focused on leveraging the generative prior of diffusion models for editing; SDEdit [\texttt{arXiv:2108.01073}] introduced Stochastic Differential Editing, which utilizes Stochastic Differential Equations (SDE) to add noise to an input and subsequently denoise it, enabling guided synthesis without task-specific training. While effective at balancing realism and faithfulness to user input [\texttt{arXiv:2108.01073}], this approach differs from the more structured multi-tasking approach seen in Palette [\texttt{arXiv:2111.05826}]. Palette provides a unified conditional diffusion framework for diverse tasks—including colorization, inpainting, and JPEG restoration—without requiring task-specific architectures or hyper-parameter tuning [\texttt{arXiv:2111.05826}].

Recent literature has sought to refine these processes by addressing specific structural and efficiency constraints. For instance, while standard diffusion models often struggle with high-resolution stability, Simple Diffusion [\texttt{arXiv:2301.11093}] achieves end-to-end high-resolution generation in pixel space by optimizing noise schedules and architecture scaling, avoiding the complexity of latent or cascaded models [\texttt{arXiv:2301.11093}]. This focus on structural integrity is furthered by the introduction of edge-preserving noise [\texttt{arXiv:2410.01540}], which employs an edge-aware scheduler to transition from edge-preserving to isotropic noise, thereby capturing fine structural details more effectively than standard Gaussian noise [\texttt{arXiv:2410.01540}].

Efficiency and sampling stability have also become central themes. To accelerate generation, research into analytical image attenuation [\texttt{arXiv:2306.13720}] proposes simultaneous image-to-zero and zero-to-noise mapping, reducing the required denoising steps [\texttt{arXiv:2306.13720}]. In the context of restoration, Diffusion Posterior Proximal Sampling (DPPS) [\texttt{arXiv:2402.16907}] improves output stability and prevents over-smoothing by using adaptive candidate selection based on the signal-to-noise ratio (SNR) and measurement identity [\texttt{arXiv:2402.16907}].

Furthermore, the field is moving toward more unified and multi-modal representations. GUD [\texttt{arXiv:2410.02667}] proposes a unified framework that utilizes component-wise noise scheduling and flexible bases (such as Fourier or wavelet) to bridge diffusion and autoregressive models [\texttt{arXiv:2410.02667}]. Complementing this, recent work on joint image-feature synthesis [\texttt{arXiv:2504.16064}] enhances generative quality by modeling low-level VAE latents alongside high-level semantic DINO features, effectively bridging representation learning with generative modeling [\texttt{arXiv:2504.16064}]. Collectively, these developments demonstrate a shift from simple stochastic editing toward highly controlled, efficient, and structurally aware conditional generation.

\subsection{Expanding Dimensions: Video and 3D Generation}

The extension of diffusion models from static images to temporal and volumetric dimensions has introduced significant complexities regarding coherence and computational efficiency. Early advancements in video generation focused on extending standard architectures through joint training on image and video data, utilizing conditional sampling for spatial and temporal extension to ensure temporal coherence [\texttt{arXiv:2204.03458}]. This trajectory evolved toward high-definition synthesis, as seen in Imagen Video, which employs a cascade of fully-convolutional spatial and temporal super-resolution models and v-parameterization to achieve high fidelity and controllability [\texttt{arXiv:2210.02303}].

A major shift in the field involves the transition from pixel-space to latent-space modeling to manage the high computational costs of video. Latent Video Diffusion Models (LVDMs) achieve high-fidelity, long-duration video (up to 1,000+ frames) by performing hierarchical diffusion within a low-dimensional 3D latent space [\texttt{arXiv:2211.13221}]. This approach is further refined in Stable Video Diffusion, which utilizes a three-stage training strategy—comprising text-to-image pretraining, video pretraining, and high-quality video finetuning—to scale LVDMs to large datasets [\texttt{arXiv:2311.15127}]. While these models prioritize scale and fidelity, other research focuses on structural control. For instance, ModelScope incorporates spatio-temporal blocks into a denoising UNet to ensure smooth movement transitions [\texttt{arXiv:2308.06571}], and other methods utilize monocular depth estimates to guide video synthesis, allowing for content editing while retaining structural integrity [\texttt{arXiv:2302.03011}].

Recent literature has also begun to isolate image-to-video (I2V) as a distinct research domain. A systematic taxonomy identifies four design pillars for I2V: condition encoding, temporal modeling, noise prior design, and spatial-temporal upsampling [\texttt{arXiv:2605.17248}]. While this categorization provides granular insight into conditioning modalities and datasets, it lacks original empirical benchmarking, leaving a "reproducibility gap" due to the prevalence of proprietary datasets [\texttt{arXiv:2605.17248}]. Furthermore, while many models focus on generation, newer architectures like Video Diffusion Transformers attempt to provide explicit control over motion dynamics by augmenting pretrained transformers with lightweight temporal modules for time editing [\texttt{arXiv:2606.10183}].

Parallel to video, the expansion into 3D generation has been achieved by repurposing 2D priors. DreamFusion demonstrates that text-to-3D synthesis can be performed without large-scale 3D datasets by using probability density distillation to optimize a Neural Radiance Field (NeRF) via a pretrained 2D diffusion model [\texttt{arXiv:2209.14988}]. This distinguishes 3D generation from video generation; whereas video models primarily struggle with temporal coherence and frame consistency, 3D methods focus on distilling 2D spatial information into a consistent volumetric representation.

\subsection{The Frontier: Manifold Geometry and Distributional Modeling}

Recent advancements in diffusion models have shifted from standard Gaussian denoising toward more sophisticated treatments of the underlying data manifold and the probabilistic nature of the reverse process. A primary tension in this frontier is the balance between sampling efficiency, distributional accuracy, and the geometric constraints of the data manifold.

One significant direction involves refining the sampling dynamics through higher-order geometric approximations. While traditional methods rely on first-order Langevin dynamics, \texttt{arXiv:2505.24222} introduces the Levenberg-Marquardt-Langevin (LML) method. This approach utilizes a low-rank approximation of the diffusion Hessian and a damping mechanism to capture second-order Hessian geometry, aiming to enhance generation quality with negligible computational overhead. This focus on the local geometry of the score function contrasts with the work of \texttt{arXiv:2606.00094}, which proposes the MIND model. Instead of optimizing the sampling trajectory via second-order dynamics, MIND explicitly models manifold geometry by integrating discrete patch tokenization into a continuous diffusion score function. By employing a soft top-k aggregation mechanism and dual-branch high-frequency feature embedding, MIND attempts to sample from a dense, low-dimensional, and compact parameterization space, effectively bridging discrete and continuous representations.

Parallel to these geometric refinements is the effort to improve the probabilistic fidelity of the reverse process. \texttt{arXiv:2502.02483} moves beyond simple mean estimation by employing scoring rules to learn the full posterior distribution of clean data samples. This distributional approach aims to accelerate inference with minimal quality degradation, representing a departure from standard regression losses. This focus on the posterior distribution is complemented by the work of \texttt{arXiv:2502.01189}, which modifies the reverse diffusion process by replacing standard Gaussian noise sampling with selections from pre-defined codebooks of fixed iid Gaussian vectors. While \texttt{arXiv:2502.02483} optimizes the \textit{learning} of the posterior, \texttt{arXiv:2502.01189} optimizes the \textit{sampling} mechanism to achieve high-quality generation paired with lossless bit-stream compression.

The intersection of diffusion and compression is further explored through stochastic differential equations (SDEs). \texttt{arXiv:2505.20984} reinterprets the compression process as a forward diffusion path, allowing for reconstruction without Gaussian noise initialization through a rate-variable feature distribution. This provides a framework for smooth rate adjustment and photo-realistic reconstruction. Finally, the relationship between the training scale and model behavior is addressed by \texttt{arXiv:2502.21278}. By utilizing "Ambient Diffusion"—training at large noise scales—this work demonstrates that memorization is primarily a requirement for denoising at low noise scales, thereby offering a method to improve the fidelity-memorization trade-off without sacrificing image quality. Together, these works suggest a transition from treating diffusion as a simple denoising task toward treating it as a complex interplay of manifold geometry, distributional estimation, and information theory.

\subsection{Tensions and Open Problems}

A central tension in the field concerns the fundamental trade-off between sample quality, training stability, and computational efficiency, often referred to as the generative learning trilemma. While early architectures like those proposed in \texttt{arXiv:1710.10196} focused on progressive growing to stabilize GAN training, and others like \texttt{arXiv:1807.03039} utilized invertible flows for exact likelihood estimation, these approaches often struggled to balance all three pillars simultaneously. Recent critiques suggest that traditional methods, including the maximum likelihood training seen in \texttt{arXiv:2101.09258} or the variational frameworks in \texttt{arXiv:2107.00630}, fail to reach the efficiency and scalability of modern diffusion-based approaches. Specifically, \texttt{arXiv:2112.07804} argues that denoising diffusion models provide a superior path to resolving this trilemma, outperforming the stability of GANs described in \texttt{arXiv:1606.03498} and the structural constraints of discrete state-space models in \texttt{arXiv:2107.03006}.

Another significant debate revolves around the optimal methodology for accelerating sampling and improving the fidelity of generative trajectories. On one side, researchers have focused on distillation techniques to reduce step counts, as seen in \texttt{arXiv:2202.00512}, or the use of specialized ODE solvers like the DPM-Solver in \texttt{arXiv:2206.00927} to navigate the reverse process more efficiently. However, these methods are often critiqued for their reliance on approximating the underlying score function. For instance, \texttt{arXiv:2209.03003} suggests that learning to "straighten" flows can offer a more direct path to generation than the complex design spaces explored in \texttt{arXiv:2206.00364}. This tension is further complicated by the emergence of flow matching; while \texttt{arXiv:2210.02747} proposes a unified framework for generative modeling, newer perspectives in \texttt{arXiv:2309.16948} suggest that denoising diffusion bridge models may offer more robust connectivity between data distributions.

Finally, there remains a growing disagreement regarding the role of guidance and latent representations in achieving photorealism. The effectiveness of classifier-free guidance, as established in \texttt{arXiv:2207.12598}, has become a standard for text-to-image synthesis, yet it is often contrasted with the heavy architectural requirements of large-scale language-guided models like those in \texttt{arXiv:2205.11487}. Furthermore, as the field moves toward more complex modalities, a tension has emerged between modeling data in high-dimensional pixel space versus compressed latent spaces. While some argue for the efficiency of latent-based diffusion, others question whether such compression sacrifices the fine-grained geometric fidelity required for truly high-resolution or manifold-aware generation.

\subsection{Outlook}

The frontier of generative modeling is rapidly shifting from static image synthesis toward highly dynamic, temporally coherent video generation and structured geometric modeling. As evidenced by recent breakthroughs in video diffusion transformers (\texttt{arXiv:2606.10183}) and the explicit modeling of data manifold geometry (\texttt{arXiv:2606.00094}), the next generation of models will likely move beyond simple denoising toward interpretable, trajectory-aware frameworks, such as hierarchical Koopman diffusion (\texttt{arXiv:2510.12220}). We expect the field to converge on a synthesis of efficient latent compression and advanced second-order sampling techniques (\texttt{arXiv:2505.24222}) to achieve real-time, high-fidelity generation across diverse, non-Euclidean data structures.

\section*{References}
\begin{enumerate}\small
  \item Auto-Encoding Variational Bayes (2013). \texttt{arXiv:1312.6114}
  \item ImageNet Large Scale Visual Recognition Challenge (2014). \texttt{arXiv:1409.0575}
  \item Microsoft COCO: Common Objects in Context (2014). \texttt{arXiv:1405.0312}
  \item Adam: A Method for Stochastic Optimization (2014). \texttt{arXiv:1412.6980}
  \item Deep Learning Face Attributes in the Wild (2014). \texttt{arXiv:1411.7766}
  \item Generative Adversarial Networks (2014). \texttt{arXiv:1406.2661}
  \item Rethinking the Inception Architecture for Computer Vision (2015). \texttt{arXiv:1512.00567}
  \item Deep Unsupervised Learning using Nonequilibrium Thermodynamics (2015). \texttt{arXiv:1503.03585}
  \item LSUN: Construction of a Large-scale Image Dataset using Deep Learning
  with Humans in the Loop (2015). \texttt{arXiv:1506.03365}
  \item U-Net: Convolutional Networks for Biomedical Image Segmentation (2015). \texttt{arXiv:1505.04597}
  \item Improved Techniques for Training GANs (2016). \texttt{arXiv:1606.03498}
  \item GANs Trained by a Two Time-Scale Update Rule Converge to a Local Nash
  Equilibrium (2017). \texttt{arXiv:1706.08500}
  \item Progressive Growing of GANs for Improved Quality, Stability, and
  Variation (2017). \texttt{arXiv:1710.10196}
  \item Attention Is All You Need (2017). \texttt{arXiv:1706.03762}
  \item Decoupled Weight Decay Regularization (2017). \texttt{arXiv:1711.05101}
  \item Diffractive Scattering: Problems in Theory and Praxis (2017). \texttt{arXiv:1711.00226}
  \item Large Scale GAN Training for High Fidelity Natural Image Synthesis (2018). \texttt{arXiv:1809.11096}
  \item A Style-Based Generator Architecture for Generative Adversarial Networks (2018). \texttt{arXiv:1812.04948}
  \item Glow: Generative Flow with Invertible 1x1 Convolutions (2018). \texttt{arXiv:1807.03039}
  \item The Unreasonable Effectiveness of Deep Features as a Perceptual Metric (2018). \texttt{arXiv:1801.03924}
  \item Neural Ordinary Differential Equations (2018). \texttt{arXiv:1806.07366}
  \item Generative Modeling by Estimating Gradients of the Data Distribution (2019). \texttt{arXiv:1907.05600}
  \item Analyzing and Improving the Image Quality of StyleGAN (2019). \texttt{arXiv:1912.04958}
  \item Sliced Score Matching: A Scalable Approach to Density and Score
  Estimation (2019). \texttt{arXiv:1905.07088}
  \item Improved Precision and Recall Metric for Assessing Generative Models (2019). \texttt{arXiv:1904.06991}
  \item Denoising Diffusion Probabilistic Models (2020). \texttt{arXiv:2006.11239}
  \item Score-Based Generative Modeling through Stochastic Differential
  Equations (2020). \texttt{arXiv:2011.13456}
  \item Denoising Diffusion Implicit Models (2020). \texttt{arXiv:2010.02502}
  \item Improved Techniques for Training Score-Based Generative Models (2020). \texttt{arXiv:2006.09011}
  \item WaveGrad: Estimating Gradients for Waveform Generation (2020). \texttt{arXiv:2009.00713}
  \item DiffWave: A Versatile Diffusion Model for Audio Synthesis (2020). \texttt{arXiv:2009.09761}
  \item Training Generative Adversarial Networks with Limited Data (2020). \texttt{arXiv:2006.06676}
  \item Taming Transformers for High-Resolution Image Synthesis (2020). \texttt{arXiv:2012.09841}
  \item Diffusion Models Beat GANs on Image Synthesis (2021). \texttt{arXiv:2105.05233}
  \item High-Resolution Image Synthesis with Latent Diffusion Models (2021). \texttt{arXiv:2112.10752}
  \item Improved Denoising Diffusion Probabilistic Models (2021). \texttt{arXiv:2102.09672}
  \item Learning Transferable Visual Models From Natural Language Supervision (2021). \texttt{arXiv:2103.00020}
  \item GLIDE: Towards Photorealistic Image Generation and Editing with
  Text-Guided Diffusion Models (2021). \texttt{arXiv:2112.10741}
  \item Cascaded Diffusion Models for High Fidelity Image Generation (2021). \texttt{arXiv:2106.15282}
  \item Image Super-Resolution via Iterative Refinement (2021). \texttt{arXiv:2104.07636}
  \item Variational Diffusion Models (2021). \texttt{arXiv:2107.00630}
  \item Zero-Shot Text-to-Image Generation (2021). \texttt{arXiv:2102.12092}
  \item SDEdit: Guided Image Synthesis and Editing with Stochastic Differential
  Equations (2021). \texttt{arXiv:2108.01073}
  \item Score-based Generative Modeling in Latent Space (2021). \texttt{arXiv:2106.05931}
  \item Knowledge Distillation in Iterative Generative Models for Improved
  Sampling Speed (2021). \texttt{arXiv:2101.02388}
  \item Palette: Image-to-Image Diffusion Models (2021). \texttt{arXiv:2111.05826}
  \item Maximum Likelihood Training of Score-Based Diffusion Models (2021). \texttt{arXiv:2101.09258}
  \item Gotta Go Fast When Generating Data with Score-Based Models (2021). \texttt{arXiv:2105.14080}
  \item ILVR: Conditioning Method for Denoising Diffusion Probabilistic Models (2021). \texttt{arXiv:2108.02938}
  \item Noise Estimation for Generative Diffusion Models (2021). \texttt{arXiv:2104.02600}
  \item On Fast Sampling of Diffusion Probabilistic Models (2021). \texttt{arXiv:2106.00132}
  \item Structured Denoising Diffusion Models in Discrete State-Spaces (2021). \texttt{arXiv:2107.03006}
  \item Tackling the Generative Learning Trilemma with Denoising Diffusion GANs (2021). \texttt{arXiv:2112.07804}
  \item Vector Quantized Diffusion Model for Text-to-Image Synthesis (2021). \texttt{arXiv:2111.14822}
  \item Learning to Efficiently Sample from Diffusion Probabilistic Models (2021). \texttt{arXiv:2106.03802}
  \item Argmax Flows and Multinomial Diffusion: Learning Categorical
  Distributions (2021). \texttt{arXiv:2102.05379}
  \item Score-Based Generative Modeling with Critically-Damped Langevin
  Diffusion (2021). \texttt{arXiv:2112.07068}
  \item Diffusion Probabilistic Models for 3D Point Cloud Generation (2021). \texttt{arXiv:2103.01458}
  \item Hierarchical Text-Conditional Image Generation with CLIP Latents (2022). \texttt{arXiv:2204.06125}
  \item Classifier-Free Diffusion Guidance (2022). \texttt{arXiv:2207.12598}
  \item Photorealistic Text-to-Image Diffusion Models with Deep Language
  Understanding (2022). \texttt{arXiv:2205.11487}
  \item Video Diffusion Models (2022). \texttt{arXiv:2204.03458}
  \item Elucidating the Design Space of Diffusion-Based Generative Models (2022). \texttt{arXiv:2206.00364}
  \item Progressive Distillation for Fast Sampling of Diffusion Models (2022). \texttt{arXiv:2202.00512}
  \item Imagen Video: High Definition Video Generation with Diffusion Models (2022). \texttt{arXiv:2210.02303}
  \item DPM-Solver: A Fast ODE Solver for Diffusion Probabilistic Model Sampling
  in Around 10 Steps (2022). \texttt{arXiv:2206.00927}
  \item Scalable Diffusion Models with Transformers (2022). \texttt{arXiv:2212.09748}
  \item Pseudo Numerical Methods for Diffusion Models on Manifolds (2022). \texttt{arXiv:2202.09778}
  \item Flow Straight and Fast: Learning to Generate and Transfer Data with
  Rectified Flow (2022). \texttt{arXiv:2209.03003}
  \item DreamBooth: Fine Tuning Text-to-Image Diffusion Models for
  Subject-Driven Generation (2022). \texttt{arXiv:2208.12242}
  \item Fast Sampling of Diffusion Models with Exponential Integrator (2022). \texttt{arXiv:2204.13902}
  \item LAION-5B: An open large-scale dataset for training next generation
  image-text models (2022). \texttt{arXiv:2210.08402}
  \item On Distillation of Guided Diffusion Models (2022). \texttt{arXiv:2210.03142}
  \item Latent Video Diffusion Models for High-Fidelity Long Video Generation (2022). \texttt{arXiv:2211.13221}
  \item Flow Matching for Generative Modeling (2022). \texttt{arXiv:2210.02747}
  \item Analytic-DPM: an Analytic Estimate of the Optimal Reverse Variance in
  Diffusion Probabilistic Models (2022). \texttt{arXiv:2201.06503}
  \item Flexible Diffusion Modeling of Long Videos (2022). \texttt{arXiv:2205.11495}
  \item New features in the differential cross sections measured at the LHC (2022). \texttt{arXiv:2210.02970}
  \item eDiff-I: Text-to-Image Diffusion Models with an Ensemble of Expert
  Denoisers (2022). \texttt{arXiv:2211.01324}
  \item DPM-Solver++: Fast Solver for Guided Sampling of Diffusion Probabilistic Models (2022). \texttt{arXiv:2211.01095}
  \item DreamFusion: Text-to-3D using 2D Diffusion (2022). \texttt{arXiv:2209.14988}
  \item Diffusion Model for Generative Image Denoising (2023). \texttt{arXiv:2302.02398}
  \item Nested Diffusion Processes for Anytime Image Generation (2023). \texttt{arXiv:2305.19066}
  \item Denoising Diffusion Bridge Models (2023). \texttt{arXiv:2309.16948}
  \item Simple diffusion: End-to-end diffusion for high resolution images (2023). \texttt{arXiv:2301.11093}
  \item Diffusion Models Generate Images Like Painters: an Analytical Theory of
  Outline First, Details Later (2023). \texttt{arXiv:2303.02490}
  \item Simultaneous Image-to-Zero and Zero-to-Noise: Diffusion Models with
  Analytical Image Attenuation (2023). \texttt{arXiv:2306.13720}
  \item Image Embedding for Denoising Generative Models (2023). \texttt{arXiv:2301.07485}
  \item Improving Denoising Diffusion Probabilistic Models via Exploiting Shared
  Representations (2023). \texttt{arXiv:2311.16353}
  \item Restoration based Generative Models (2023). \texttt{arXiv:2303.05456}
  \item Align your Latents: High-Resolution Video Synthesis with Latent
  Diffusion Models (2023). \texttt{arXiv:2304.08818}
  \item Adding Conditional Control to Text-to-Image Diffusion Models (2023). \texttt{arXiv:2302.05543}
  \item Structure and Content-Guided Video Synthesis with Diffusion Models (2023). \texttt{arXiv:2302.03011}
  \item ModelScope Text-to-Video Technical Report (2023). \texttt{arXiv:2308.06571}
  \item SDXL: Improving Latent Diffusion Models for High-Resolution Image
  Synthesis (2023). \texttt{arXiv:2307.01952}
  \item Stable Video Diffusion: Scaling Latent Video Diffusion Models to Large
  Datasets (2023). \texttt{arXiv:2311.15127}
  \item About AKM scaling and oscillations in elastic scattering at very small
  momentum transfer at the LHC (2023). \texttt{arXiv:2307.15445}
  \item Accelerated Image-Aware Generative Diffusion Modeling (2024). \texttt{arXiv:2408.08306}
  \item Blue noise for diffusion models (2024). \texttt{arXiv:2402.04930}
  \item Optical Diffusion Models for Image Generation (2024). \texttt{arXiv:2407.10897}
  \item DiffuseHigh: Training-free Progressive High-Resolution Image Synthesis
  through Structure Guidance (2024). \texttt{arXiv:2406.18459}
  \item Enhancing Diffusion Models for High-Quality Image Generation (2024). \texttt{arXiv:2412.14422}
  \item GUD: Generation with Unified Diffusion (2024). \texttt{arXiv:2410.02667}
  \item Generative AI in Vision: A Survey on Models, Metrics and Applications (2024). \texttt{arXiv:2402.16369}
  \item Edge-preserving noise for diffusion models (2024). \texttt{arXiv:2410.01540}
  \item Diffusion Posterior Proximal Sampling for Image Restoration (2024). \texttt{arXiv:2402.16907}
  \item Generative Modeling with Diffusion (2024). \texttt{arXiv:2412.10948}
  \item Distributional Diffusion Models with Scoring Rules (2025). \texttt{arXiv:2502.02483}
  \item Compressed Image Generation with Denoising Diffusion Codebook Models (2025). \texttt{arXiv:2502.01189}
  \item Generative Image Compression by Estimating Gradients of the Rate-variable Feature Distribution (2025). \texttt{arXiv:2505.20984}
  \item Boosting Generative Image Modeling via Joint Image-Feature Synthesis (2025). \texttt{arXiv:2504.16064}
  \item Unleashing High-Quality Image Generation in Diffusion Sampling Using Second-Order Levenberg-Marquardt-Langevin (2025). \texttt{arXiv:2505.24222}
  \item MSF: Efficient Diffusion Model Via Multi-Scale Latent Factorize (2025). \texttt{arXiv:2501.13349}
  \item Hierarchical Koopman Diffusion: Fast Generation with Interpretable Diffusion Trajectory (2025). \texttt{arXiv:2510.12220}
  \item Redacted (2025). \texttt{arXiv:2506.17293}
  \item Does Generation Require Memorization? Creative Diffusion Models using Ambient Diffusion (2025). \texttt{arXiv:2502.21278}
  \item Denoising Score Distillation: From Noisy Diffusion Pretraining to One-Step High-Quality Generation (2025). \texttt{arXiv:2503.07578}
  \item Diffusion Image Generation with Explicit Modeling of Data Manifold Geometry (2026). \texttt{arXiv:2606.00094}
  \item Image-to-Video Diffusion: From Foundations to Open Frontiers (2026). \texttt{arXiv:2605.17248}
  \item Making Time Editable in Video Diffusion Transformers (2026). \texttt{arXiv:2606.10183}
\end{enumerate}

\end{document}
